% from ega0-22-fr.tex
\oldpage[0\textsubscript{IV-1}]{184}

(ii) Si $B$ est une $\Lambda$-algèbre formellement lisse pour la topologie
$\mathfrak J$-préadique, et $A$ une $\Lambda$-algèbre formellement lisse pour
la topologie discrète, $B$ est une $A$-algèbre formellement lisse pour la
topologie $\mathfrak J$-préadique (19.3.5, (ii))~; les conditions $\alpha$) et
$\beta$) résultent donc de (19.5.6). En outre, (20.7.2) montre que
$u_{B/A/\Lambda}$ est formellement inversible à gauche, donc
$u_{B/A/\Lambda}\otimes 1_C$ est inversible à gauche (19.1.7), et \emph{a
fortiori} injectif. Enfin, $\Omega^1_{B/A}$ est un $B$-module formellement
projectif (20.4.9), et par suite $\Omega^1_{B/A}\otimes_BC$ est un $C$-module
projectif (19.2.4).
\end{proof}

\begin{corollary}[22.1.2]
\label{0.22.1.2-fr}
Soient $s:\Lambda\to A$, $u:A\to B$ deux homomorphismes d'anneaux,
$\mathfrak m$ un idéal maximal de $B$, $K=B/\mathfrak m$ le corps quotient.
On suppose que les $\Lambda$-algèbres $A$ et $K$ soient formellement lisses
pour la topologie discrète. Alors les conditions suivantes sont
équivalentes~:
\begin{enumerate}
  \item[a)] $B$ est une $A$-algèbre formellement lisse pour la topologie
  $\mathfrak m$-préadique.
  \item[b)] L'homomorphisme canonique
  \begin{equation}
  \tag{22.1.2.1}
  \label{0.22.1.2.1-fr}
    \mathbf S_K^*(\mathfrak m/\mathfrak m^2)
    \longrightarrow\operatorname{gr}_{\mathfrak m}^*(B)
  \end{equation}
  est bijectif, et l'homomorphisme caractéristique
  \begin{equation}
  \tag{22.1.2.2}
  \label{0.22.1.2.2-fr}
    \chi:\Upsilon_{K/B/\Lambda}\longrightarrow\mathfrak m/\mathfrak m^2
  \end{equation}
  est injectif.
\end{enumerate}
\end{corollary}

\begin{proof}
Cela résulte aussitôt de (22.1.1), les conditions $\alpha$) et $\delta$)
étant ici trivialement vérifiées, puisque $K$ est un corps.
\end{proof}

\begin{remark}[22.1.3]
\label{0.22.1.3-fr}
Supposons vérifiées les hypothèses de (22.1.2) et en outre supposons que
$\mathfrak m/\mathfrak m^2$ soit un $K$-espace vectoriel de rang \emph{fini}.
Alors, pour que $B$ soit une $A$-algèbre formellement lisse pour la topologie
$\mathfrak m$-préadique, il suffit que les conditions suivantes soient
satisfaites~:

\begin{env}[22.1.3.1]
\label{0.22.1.3.1-fr}
Étant donnés une algèbre de polynômes $E=K[T_1,\ldots,T_n]$, l'idéal maximal
$\mathfrak n$ de $E$ engendré par les $T_i$ $(1\leqslant i\leqslant n)$, et
un idéal $\mathfrak b\supset\mathfrak n^k$ dans $E$, tout
$\Lambda$-homomorphisme $v:B\to E/\mathfrak b$ qui, par passage aux
quotients, donne l'identité $K\to K$, se factorise en
\[
  B\xrightarrow{w}E/\mathfrak n^k\longrightarrow E/\mathfrak b,
\]
où $w$ est un $\Lambda$-homomorphisme.
\end{env}

\begin{env}[22.1.3.2]
\label{0.22.1.3.2-fr}
Si $F$ est l'extension triviale $D_K(K)$ («~algèbre des nombres duaux sur
$K$~»), $\rho:A\to F$ un homomorphisme d'anneaux tel que le composé
$A\xrightarrow{\rho}F\to K$ soit l'homomorphisme canonique, alors tout
$\Lambda$-homomorphisme $v:B\to K$ qui, par passage aux quotients, donne
l'identité $K\to K$, se factorise en
\[
  B\xrightarrow{w}F\longrightarrow K,
\]
où $w$ est un $\Lambda$-homomorphisme (pour la structure de $A$-algèbre sur
$F$ définie par $\rho$).
\end{env}

En effet, on a vu (19.5.5) que la condition (22.1.3.1) entraîne que
l'homomorphisme canonique (22.1.2.1) est bijectif. En vertu de (22.1.2), il
suffit ensuite de voir que l'homomorphisme canonique
\[
  u_{B/A/\Lambda}\otimes1_K:
  \Omega^1_{A/\Lambda}\otimes_AK
  \longrightarrow\Omega^1_{B/\Lambda}\otimes_BK
\]
est injectif, ou encore (puisqu'il s'agit d'espaces vectoriels sur $K$) que
pour tout $K$-espace vectoriel $L$, l'homomorphisme canonique
$\operatorname{D\acute{e}r}_\Lambda(B,L)\to\operatorname{D\acute{e}r}_\Lambda(A,L)$ est
\emph{surjectif} (en vertu de (20.4.8))~; il est clair que pour cela, il faut
et il suffit que l'homomorphisme canonique
$\operatorname{D\acute{e}r}_\Lambda(B,K)\to\operatorname{D\acute{e}r}_\Lambda(A,K)$ soit
surjectif, d'où notre assertion, en vertu de (20.1.5).

On notera que les deux conditions (22.1.3.1) et (22.1.3.2) sont entraînées
par la suivante~:

\begin{env}[22.1.3.3]
\label{0.22.1.3.3-fr}
Pour toute $\Lambda$-algèbre locale artinienne $E$ de corps résiduel égal à
$K$, et tout idéal nilpotent $\mathfrak r$ de $E$, tout
$\Lambda$-homomorphisme $B\to E/\mathfrak r$ qui, par passage aux quotients,
donne l'identité $K\to K$, se factorise en
$B\xrightarrow{u}E\to E/\mathfrak r$, où $u$ est un
$\Lambda$-homomorphisme.
\end{env}
\end{remark}

